\chapter{从裸机到可登录}
\label{ch:server-access}
\chapterlead{服务器的第一件事是把门做窄：只开必要的端口，只允许密钥登录，只给必要的权限。本章在一台新装服务器上完成这些设置，并说明每一项挡住了什么。}

\section{起步：新服务器的五项检查}

\begin{shellsession}
$ ssh weiwen@203.0.113.10
$ nixos-version
$ ip -brief address
$ ss -tlnp
$ sudo journalctl -u sshd -n 20
\end{shellsession}

五条命令确认：能登录、版本正确、地址可达、有哪些端口在监听、SSH 服务是否正常。把这五项做成脚本，每次接手新机器先跑一遍。

\section{SSH 加固}

\begin{nixcode}
services.openssh = {
  enable = true;
  openFirewall = true;
  settings = {
    PasswordAuthentication = false;
    KbdInteractiveAuthentication = false;
    PermitRootLogin = "no";
    MaxAuthTries = 3;
    AllowUsers = [ "weiwen" ];
  };
};
\end{nixcode}

四项设置各挡住一类攻击。关闭密码登录，使暴力猜测无效，只有持有私钥的人能连；禁止 root 直接登录，使攻击者必须先猜中普通用户名；\texttt{MaxAuthTries} 限制每次连接的重试次数；\texttt{AllowUsers} 把可登录账户收敛到白名单。

\begin{warnbox}[改 SSH 配置的操作顺序]
保持当前会话不要退出。另开一个终端，用新配置连接；确认成功后再关闭旧会话。若新连接失败，旧会话可以改回配置并重建。这一步能避免把自己锁在服务器外面。
\end{warnbox}

\section{公钥的三种放法}

\begin{nixcode}
# 一、写进系统配置，重建后生效，重装也还在
users.users.weiwen.openssh.authorizedKeys.keys = [
  "ssh-ed25519 AAAA... weiwen@laptop"
  "ssh-ed25519 AAAA... weiwen@desktop"
];

# 二、临时追加，立即生效，重建后可能被覆盖
# $ echo "ssh-ed25519 ..." >> ~/.ssh/authorized_keys

# 三、从网络拉取（NixOS 特性）
users.users.weiwen.openssh.authorizedKeys.keys =
  [ (builtins.readFile ./keys/weiwen.pub) ];
\end{nixcode}

第一种最稳：公钥随配置走，重建与重装都不会丢。第二种适合临时给同事开权限。第三种把公钥文件放进仓库，便于多人维护。

\section{防火墙}

\begin{nixcode}
networking.firewall = {
  enable = true;
  allowedTCPPorts = [ 22 80 443 ];
  # 只对特定网段开放
  extraCommands = ''
    iptables -A nixos-fw -p tcp -s 10.0.0.0/8 --dport 5432 -j nixos-fw-accept
  '';
};
\end{nixcode}

默认拒绝入站，显式列出允许的端口。数据库这类只应内网访问的服务，不要直接加进 \texttt{allowedTCPPorts}，而是按来源网段放行，或根本不对外开放、通过 SSH 隧道访问。

\begin{shellsession}
$ sudo nft list ruleset | head -20
$ ss -tlnp
\end{shellsession}

改完防火墙后用 \texttt{ss} 确认监听，用 \texttt{nft} 确认规则。两者不一致时，先看服务是否真的监听了外部地址。

\section{用户与权限}

\begin{nixcode}
users.users.deploy = {
  isNormalUser = true;
  group = "deploy";
  openssh.authorizedKeys.keys = [ "ssh-ed25519 AAAA... ci@runner" ];
  # 只用于部署，不给 sudo
};

users.groups.deploy = {};

security.sudo.extraRules = [{
  users = [ "weiwen" ];
  commands = [
    { command = "/run/current-system/sw/bin/systemctl restart web"; options = [ "NOPASSWD" ]; }
  ];
}];
\end{nixcode}

原则是最小权限：部署账户只拥有部署目录的写权限，不进入 \texttt{wheel} 组；需要频繁执行的单条命令可以单独授权，免去每次输入密码，但授权范围要写到具体命令。

\section{时间、时区与主机名}

\begin{nixcode}
time.timeZone = "Asia/Shanghai";
services.chrony.enable = true;
networking.hostName = "srv1";
networking.domain = "example.com";
\end{nixcode}

时间同步看起来琐碎，但日志排序、证书校验、备份时间戳都依赖它。\texttt{chrony} 默认启用，确认它处于运行状态即可。

\begin{shellsession}
$ timedatectl
$ hostnamectl
\end{shellsession}

\section{练习}

\begin{enumerate}
  \item 在一台测试机上关闭密码登录，用两个终端验证“保留旧会话”的流程。
  \item 用 \texttt{ss -tlnp} 列出所有监听端口，逐个判断是否应当对外开放。
  \item 为部署账户写配置，使它能写 \texttt{/srv/www} 但不能 sudo。
  \item 用 \texttt{nft list ruleset} 查看默认规则，说明没有显式放行的端口会发生什么。
  \item 用 \texttt{timedatectl} 确认时间同步状态，手动改一次时区并观察输出变化。
\end{enumerate}

\begin{checkpoint}
本章结束时，应能：
\begin{itemize}
  \item 用五条命令完成新服务器的初次检查；
  \item 写出加固后的 SSH 配置，说明每项挡住了什么；
  \item 用三种方式管理公钥，判断何时用哪一种；
  \item 配置防火墙，按来源而非仅按端口放行；
  \item 按最小权限原则创建账户并限制 sudo 范围。
\end{itemize}
\end{checkpoint}
